% Original r6 appendices, unchanged. Integrate into the corrected companion before release.
\appendix
\section{Excluding a constant negative magnitude}\label{sec:constant}

Suppose the error is negative with a constant magnitude, $E_n=-m$ for a fixed
$m>0$ and every $n$.  By
Proposition~\ref{res:update} the state then climbs by exactly $m$ each step,
so $C_n=c+nm$ with $c=C_0$, and the definition of the centred state becomes a
\emph{shape equation}
\[
 D_n+m=(a_n-1)\,(c+nm) .
\tag{5.1}\label{eq:shape}
\]
Together with $D_{n+1}=a_nD_n$ this is a closed system in $(a,D)$, and it has
no solutions.  The following instance shows how the failure happens; the
theorem then shows that it cannot be avoided.

\begin{example}[the shape equation running until it fails]\label{ex:shape}
Take $m=c=1$, so that $C_n=1+n$ and the shape equation~\eqref{eq:shape} reads
$D_n+1=(a_n-1)(1+n)$, and take $D_0=1$.  Each step is now determined: at $n=0$
the equation reads $2=(a_0-1)\cdot1$, so $a_0=3$, and $D_1=a_0D_0=3$; at $n=1$
it reads $4=(a_1-1)\cdot2$, so $a_1=3$, and $D_2=a_1D_1=9$; at $n=2$ it reads
$10=(a_2-1)\cdot3$, which has no integer solution, and the orbit stops.
Longer prefixes occur for other starting values: the finite search recorded in
Appendix~\ref{app:residue} found none longer than $17$ among the initial
states below $5000$.
\end{example}

\begin{theorem}[no constant negative magnitude]\label{res:constant}
There is no pair of sequences $a,D:\N\to\N$ with $a_n\ge2$ for all $n$
satisfying $D_{n+1}=a_nD_n$ and~\eqref{eq:shape}, for any $m>0$ and any
$c$.  The same holds if the shape only begins at some index.
\end{theorem}

\begin{proof}[Proof when $c$ and $m$ are coprime]
Assume first $\gcd(c,m)=1$.  Two facts about the multipliers do all the work,
and the difficulty lies in the fact that they concern the same finite set of
primes from opposite directions: the first makes every multiplier meet that
set, the second lets each prime of the set meet at most one multiplier.

\emph{Every $a_j$ shares a prime with $m$.}  Suppose $\gcd(a_j,m)=1$.  Then
$m$ is invertible modulo $a_j$, so some $n$ has $c+nm\equiv0$, that is
$a_j\mid C_n$; and $a_j\mid D_n$ for every $n>j$, since $D$ is multiplicative
with $a_j$ among its factors.  Choosing such an $n$ beyond $j$
and reading~\eqref{eq:shape} modulo $a_j$ gives $a_j\mid m$, contradicting
$\gcd(a_j,m)=1$ and $a_j\ge2$.

\emph{A prime divisor of $m$ occurs in at most one $a_j$.}  Suppose $p\mid m$
and $p\mid a_j$.  For $n>j$ we have $p\mid D_n$, so~\eqref{eq:shape} gives
$(a_n-1)C_n\equiv0 \pmod p$.  Now $C_n=c+nm\equiv c$, and $p\nmid c$ because
$p\mid m$ and $\gcd(c,m)=1$; hence $p\mid a_n-1$, so $p\nmid a_n$.  Thus $p$
cannot divide any later multiplier.

The two facts are incompatible.  Each of the infinitely many multipliers needs
a prime divisor of $m$, and each prime divisor of $m$ serves at most one
multiplier, while $m$ has only finitely many prime divisors.
\end{proof}

\noindent The case $\gcd(c,m)=1$ is the
\lword{Erdos243/ReciprocalTailRigidity.lean}{147}{no_scalePrimitiveConstantNegative_orbit}{scale-primitive
exclusion}.  The special case $m=c=1$, worked through in
Example~\ref{ex:shape}, where $m$ has no prime divisors at all and the first
fact is immediately contradictory, is recorded separately as the
\lword{Erdos243/ReciprocalTailRigidity.lean}{100}{no_normalizedConstantNegative_orbit}{normalised
exclusion}.

\begin{proof}[Removing the scale]
For general $c$, put $g=\gcd(c,m)$.  Equation~\eqref{eq:shape} shows
$g\mid D_n$ for every $n$, so dividing $D$, $c$ and $m$ by $g$ leaves a system
of the same shape with coprime data, which the previous case excludes.
\end{proof}

\noindent Formalised as the
\lword{Erdos243/ReciprocalTailRigidity.lean}{286}{no_constantNegative_orbit}{exclusion
at every scale}, and the eventual form, obtained by shifting the orbit to the
first constant index, as the
\lword{Erdos243/ReciprocalTailRigidity.lean}{350}{no_eventuallyConstantNegative_orbit}{eventual
exclusion}.

Theorem~\ref{res:constant} excludes a centred state that is eventually constant
and negative, at every magnitude and every scale.  It says nothing about a
negative state whose magnitude changes, and the argument genuinely uses
constancy: the finiteness of the set of prime divisors of $m$ is what the
pigeonhole uses, and a varying magnitude presents a fresh set at each index.

\section{Excluding a periodic negative magnitude}\label{sec:periodic}

The next case allows the magnitude to vary, provided it repeats.  Write
$e_n=-E_n>0$ for the magnitude, as in Section~\ref{sec:state}, so that
$C_{n+1}=C_n+e_n$; suppose $e$ has period $h$, meaning $e_{n+h}=e_n$ for every
$n$, and that the state gains a fixed \emph{drift} $M>0$ over one period,
meaning $C_{n+h}=C_n+M$ for every $n$.

At $h=1$ this says that the magnitude is constant and that $M$ is its value,
which is the situation of Section~\ref{sec:constant};
Theorem~\ref{res:constant} already excludes it, and does so without the extra
hypothesis $e_n<a_n$ imposed below.  The new content is $h\ge2$.

Section~\ref{sec:constant} could run its pigeonhole on the prime divisors of
the single number $m$.  Here the magnitude changes inside a period, and the
number available to run it on is the drift $M$.  Three lemmas replace the two
facts of Section~\ref{sec:constant}.  The first is
that multipliers persist: every $a_j$ divides every strictly later denominator
state, the
\lword{Erdos243/ReciprocalTailRigidity.lean}{382}{base_dvd_denState_later}{persistence
of a multiplier}.  The second is a \emph{lock}: a prime already dividing $D$ and also
dividing the current multiplier is forced into the next tail state, the
\lword{Erdos243/ReciprocalTailRigidity.lean}{402}{divisorLocksTailSucc}{one-step
lock}, and once present in both $D$ and $C$ it stays in every later $D$, $C$
and magnitude, the
\lword{Erdos243/ReciprocalTailRigidity.lean}{429}{divisorLock_persists}{persistence
of a lock}.  The third transports a lock backwards through the period: if a
divisor of the drift $M$ occurs in two multipliers, it is a common divisor of
$C_0$ and of every magnitude, the
\lword{Erdos243/ReciprocalTailRigidity.lean}{494}{repeatedDivisor_forces_periodicCommonScale}{repeated-divisor
transport}, using that both the period relation and the drift iterate over any
number of periods, the
\lword{Erdos243/ReciprocalTailRigidity.lean}{461}{periodic_add_mul}{iterated
period} and the
\lword{Erdos243/ReciprocalTailRigidity.lean}{476}{phaseDrift_add_mul}{iterated
drift}.

\begin{theorem}[no periodic negative magnitude]\label{res:periodic}
Let $a,D,C,e:\N\to\N$ with $a_n\ge2$, $e_n>0$ and $e_n<a_n$ for every $n$,
satisfying
\[
 D_{n+1}=a_nD_n,\qquad C_{n+1}=C_n+e_n,\qquad D_n+e_n=(a_n-1)C_n ,
\]
and suppose $e_{n+h}=e_n$ and $C_{n+h}=C_n+M$ for some $h>0$ and $M>0$.  This
is impossible.
\end{theorem}

\begin{proof}
Strong induction on the drift $M$.  The difficulty lies in the fact that a
repeated prime divisor of $M$ need not contradict anything directly.  It
instead forces a common scale on the whole orbit, which can be divided out, and
the induction runs on the drift that the division reduces.

Suppose first that no prime divisor of $M$ is a common divisor of $C_0$ and of
every magnitude.  Repeated-divisor transport then applies in the
contrapositive: a prime divisor of $M$ occurring in two of the multipliers
would be exactly such a common divisor, so each prime divisor of $M$ occurs in
at most one multiplier.  Against this, the argument of
Section~\ref{sec:constant} run phase by phase (along a fixed residue class
modulo $h$ the magnitude is constant and $C$ advances by $M$ at each period, so
the multipliers of that phase meet an arithmetic progression exactly as in
Section~\ref{sec:constant}) makes every multiplier share a prime with $M$.
Since $M$ has only finitely many prime divisors and there are infinitely many
multipliers, pigeonhole gives a contradiction.

Otherwise some prime $p$ divides $M$, divides $C_0$, and divides every
magnitude.  Then $p\mid C_n$ for every $n$, since $C_{n+1}=C_n+e_n$ and both
summands on the right are divisible by $p$, and the shape equation
$D_n+e_n=(a_n-1)C_n$ then gives $p\mid D_n$ for every $n$.  Divide $D$, $C$ and
$e$ by $p$.  The multipliers are untouched, so $a_n\ge2$ and $e_n<a_n$ persist
and the magnitudes stay positive; the three recurrences are homogeneous in
$(D,C,e)$ and so survive; the period is still $h$; and the drift becomes
$M/p<M$.  The inductive hypothesis applies.
\end{proof}

\noindent The first of the two cases above is the
\lword{Erdos243/ReciprocalTailRigidity.lean}{548}{no_phasePrimitivePeriodicNegative_orbit}{phase-primitive
exclusion}, the induction is the
\lword{Erdos243/ReciprocalTailRigidity.lean}{690}{no_periodicNegative_orbit}{exclusion
at every scale}, and the eventual form is the
\lword{Erdos243/ReciprocalTailRigidity.lean}{796}{no_eventuallyPeriodicNegative_orbit}{eventual
exclusion}.

The bound $e_n<a_n$ is a genuine restriction and not a normalisation.  It is
the range in which a multiplier cannot divide its own phase magnitude, which
is what the lock argument needs.  It holds in the intended reading, where
$e_n$ is a centred representative and $a_n$ grows quadratically, but it is
assumed here and not derived.  The drift hypothesis $M>0$ is also used: a
periodic magnitude with zero drift would be identically zero, which is the case
already settled by Section~\ref{sec:descent}.

\begin{remark}
On the canonical orbit of~\cite{koizumi2025}, read through the dictionary of
Section~\ref{sec:problem}, the bound $e_n<a_n$ holds eventually.  Lemma~4(2)
there gives the centring range $-C_n/2\le E_n<C_n/2$, so $e_n=|E_n|\le C_n/2$ at
a negative index, and the proof of Lemma~4(3) gives $C_{n+1}\le\tfrac32C_n$
~\cite[p.~11]{koizumi2025},
hence the uniform bound $C_n\le C_0(3/2)^{n}$.  Against this, summability of
$\sum1/a_n$ forces
$a_n\to\infty$, and the growth hypothesis $a_n^{2}/a_{n+1}\to1$ gives
$a_{n+1}\ge a_n^{2}/2$ for all large $n$, hence $a_{n+1}\ge2a_n$ once
$a_n\ge4$; so the multipliers grow at least geometrically with ratio $2$ and
outgrow the bound on $C_n$, making $e_n/a_n$ eventually small.  Two caveats.
Summability is used for this step, so the estimate is not available from the
state recurrence alone; and the conclusion is eventual and not universal,
whereas Theorem~\ref{res:periodic} imposes $e_n<a_n$ at every index, so it is
the eventual form of the exclusion cited above that applies.  Under the
periodicity hypothesis of Theorem~\ref{res:periodic} the estimate is in any
case unnecessary, since a periodic magnitude is bounded while the multipliers
tend to infinity.
\end{remark}


\section{Two-step cancellation and feedback}\label{app:feedback}

The one-step state equations hide a second-order constraint once the raw
numerator is reduced at two consecutive stages.  Write $u,u_1,u_2$ for the
primitive numerator coordinates, $v,v_1$ for the corresponding denominator
coordinates, and $h,h_1$ for the two common factors removed in reduction.  If
the two steps have multipliers $a,a_1$, then the exact hypotheses are
\[
 hu_1+v=au,\qquad hv_1=av,\qquad
 h_1u_2+v_1=a_1u_1.
\]

\begin{theorem}[two-step primitive recurrence]\label{res:secondorder}
Under these three equations,
\[
 a^2u+hh_1u_2=h(a+a_1)u_1.
\]
\end{theorem}

\begin{proof}
Multiply the first equation by $a$, replace $av$ using the second equation,
and then replace $h_1u_2+v_1$ using the third.  The remaining terms factor as
the displayed right-hand side.
\end{proof}

\noindent This is the
\lword{Erdos243/DynamicCancellation.lean}{268}{reducedStep_secondOrder}{two-step
primitive recurrence}.  The hard point is keeping both changing cancellation
factors visible.  The result is stronger than a one-step congruence, though it
remains non-autonomous and cannot by itself force $a_1=a^2-a+1$.

There is a complementary integer identity when the relevant step is
cancellation-free.  Put $q=ap-p_1$, and suppose the next primitive coordinate
satisfies
\[
 p_2+a^2p=(a+a_1)p_1.
\]
Then
\begin{theorem}[cancellation-free curvature identity]\label{res:curvature}
\[
 q^2+\bigl(pp_2-p_1^2\bigr)=(a_1-a)pp_1.
\]
\end{theorem}

The square $q^2$ is the local remainder, while $pp_2-p_1^2$ is a discrete
curvature term.  The right side records the change in multiplier, so the
identity supplies a diagnostic for an anti-shadowing argument without
asserting a sign or a global monotonicity.  It is checked as
\lword{Erdos243/DynamicCancellation.lean}{307}{cancellationFree_curvature_square}{the
cancellation-free curvature square}; no global anti-shadowing conclusion is
claimed.

Finally, the reduced denominator transport has a precise realizability gate.
Let $e,e_1$ be centred errors and let $\Delta$ be the Sylvester defect.  If
\[
 hu_1=u-e,
 \qquad he_1=a^2e-\Delta(u-e),
\]
then the transport equality
\[
 h\bigl((a_1-1)u_1+e_1\bigr)
   =a\bigl((a-1)u+e\bigr)
\]
has mismatch which factors through $u-e$.  Consequently, when $u-e\ne0$,
\[
 h\bigl((a_1-1)u_1+e_1\bigr)
   =a\bigl((a-1)u+e\bigr)
 \quad\Longleftrightarrow\quad
 a_1=a^2-a+1+\Delta.
\]
This is the exact
\lword{Erdos243/FeedbackRealizability.lean}{64}{feedback_denominator_transport_iff}{feedback
transport equivalence}, obtained from the factored defect at
\lword{Erdos243/FeedbackRealizability.lean}{15}{feedback_transport_defect}{feedback
transport defect}.  The nonzero-raw-numerator hypothesis is essential: if
$u-e=0$, the factorised mismatch cannot identify $a_1$.

These three identities require orbit-level hypotheses that remain unproved:
eventual absence of cancellation, a fixed sign for the curvature, or cofinal
passage through the feedback gate.  The global mixed-sign and
unbounded-negative regimes of Problem~\#243 remain open.


\section{Record increments, protected epochs and record-amplified errors}\label{sec:records}

This section works in the reduced coordinates of the tail and allows an
arbitrary cancellation factor at every step.  Write
$x_n=\sum_{k\ge n}1/a_k=u_n/v_n$ in lowest terms, put
$e_n=v_n-(a_n-1)u_n$ for the centred error, $w_n=a_nu_n-v_n=u_n-e_n$ for the
raw next numerator, and $h_n=\gcd(w_n,a_n^2)$ for the content removed, so that
\[
 h_nu_{n+1}=w_n,\qquad h_nv_{n+1}=a_nv_n,\qquad \gcd(u_n,u_{n+1})=1 .
\]
Under the growth hypothesis the selector is eventually the nearest integer and
$e_n/u_n\to0$, so $2|e_n|<u_n$ on a tail.  Let $R_n=\max_{k\le n}u_k$ be the
running maximum.  The \emph{true record increment} is $R_{n+1}-R_n$; the
\emph{record-setting jump} $u_{n+1}-u_n=(R_{n+1}-R_n)+(R_n-u_n)$ also recovers
the drawdown $R_n-u_n$.  The exact orbit
$(8,177)\to(7,4071)\to(10,2373393)$ with multipliers $23$ and $583$ has
records $8$ and $10$: the increment is $2$, the jump is $3$, and
$\gcd(8,10)=2$.  The distinction governs which parity arguments are available.

The results below are the finite arithmetic cores of the second return batch of
2026-09-05, together with the global criteria they support.  Four Lean modules
carry the finite cores and are listed in Appendix~\ref{app:index}; the global
criteria are ordinary proofs recorded in the working note
\texttt{WeightedRecordOctupleCriteria.md}.  Every result is conditional on its
stated hypotheses and none settles Problem~\#243.

\begin{theorem}[historical square payment]\label{res:squarepayment}
On every step and for every prime $p$,
$\nu_p(h_n)\ge2\,(\nu_p(v_n)-\nu_p(v_{n+1}))_+$.  Consequently, if $p^b\mid v_T$
and $p\nmid v_s$ for some $s>T$, then $p^{2b}$ divides $\prod_{T\le n<s}h_n$, and
this divisibility persists in every later product even if $p$ re-enters the
denominator.
\end{theorem}

\begin{proof}
Since $h_n\mid w_n$ and $h_n\mid a_nv_n$, the identity
$\gcd(w_n,a_nv_n)=\gcd(w_n,a_n^2)$ gives $h_n\mid a_n^2$, hence
$\nu_p(h_n)\le2\nu_p(a_n)$, while the denominator update gives
$\nu_p(v_{n+1})=\nu_p(a_n)+\nu_p(v_n)-\nu_p(h_n)$.  A drop of size $\delta$
therefore has $\nu_p(h_n)=\nu_p(a_n)+\delta\ge2\delta$.  Summing over the
interval telescopes the drops.
\end{proof}

The per-step inequality and the telescoped product are checked in Lean.

\begin{theorem}[unit record increments force the recurrence]\label{res:unitincrement}
Assume the hypotheses of Problem~\ref{res:problem}.  If $R_{n+1}-R_n\le1$ for
all sufficiently large $n$, then $a_{n+1}=a_n^2-a_n+1$ eventually.  Hence the
sequence is eventually Sylvester if and only if
$\#\{n:R_{n+1}-R_n\ge2\}<\infty$.  No hypothesis is imposed on drawdowns, on
record-setting jumps, or on the cancellation factors.
\end{theorem}

\begin{proof}[Proof sketch]
If the recurrence fails then $e_n\neq0$ eventually and $u_n\to\infty$.  A late
multiplier coprime to all earlier ones carries a prime power $P=p^e\mid v_s$
with $P\ge R_s+3$, because otherwise $a_j$ divides $\operatorname{lcm}(1,\ldots,R_{j+1}+2)$
and $\log a_j\le\exp(o(j))$, against $\log a_j\ge c2^j$.  Let $H$ be the least
multiple of $p$ above $R_s$, so $R_s<H\le R_s+p$.  At the first $t>s$ with
$u_t\ge H$ one has $R_{t-1}\le H-1$ and $u_t=R_t\le R_{t-1}+1\le H$, so $u_t=H$.
Before $t$ the raw numerators satisfy $w_n<\tfrac32u_n<\tfrac32H<pP$, so the
valuation threshold keeps $P\mid v_t$, and $p\mid u_t$ contradicts
$\gcd(u_t,v_t)=1$.  The converse is $u_n=1$ on a Sylvester tail.
\end{proof}

The finite cut and the composed consumer are checked in Lean with the
prime-power supply as an explicit hypothesis; the supply is an ordinary
argument in the working note, not a Lean theorem, and is not used here as a
global criterion.  The two-unit record-jump theorem of the first batch, which bounds
the record-setting jump by $2$ and needs an odd prime and a parity argument, is
incomparable with Theorem~\ref{res:unitincrement}: a jump of at most $2$ bounds
the increment by $2$ and an increment of at most $1$ allows arbitrarily large
jumps after a drawdown.  True increments of size exactly $2$ are covered by
neither theorem.

\begin{theorem}[record increments against cancellation]\label{res:kgamma}
Put $\ell(x)=\log_2\log_2\max(4,x)$, $K=\limsup(R_{n+1}-R_n)/\ell(R_n)$ and
$\Gamma=\liminf\ell(G_n)/\ell(R_n)$, where $G_n=\gcd(C_n,D_n)$.  On a
non-Sylvester orbit $K+\Gamma\ge1$.  In particular
$R_{n+1}-R_n=o(\ell(R_n))$ forces $G_n\ge2^{(\log_2R_n)^{1-\eta}}$ eventually
for every $\eta>0$.
\end{theorem}

The proof charges the retirement of fresh multipliers whose primes all leave the
reduced denominator to $G_n$ through Theorem~\ref{res:squarepayment}, and forces
a large increment by a Chinese-remainder block when they do not retire.  It is a
necessary condition compatible with $\log G_n=o(n)$ and $\log R_n=o(n)$.

\begin{theorem}[protected epoch energy]\label{res:epoch}
Let $p\ge3$ be prime, $Q=p^{\ell}\mid v_s$, $Q\ge16$, $L=pQ/2$, $R_s<L/2$, and
let $\tau>s$ satisfy $u_\tau\ge L$.  Call a step in $[s,\tau)$ charged if it is
the first to cross an odd multiple of $p$ in $(L/2,L]$; let $K$ be the number of
charged steps and $X$ the sum over them of $u_{n+1}-u_n-2$.  Every charged step
is a global record with $h_n=1$ and $u_{n+1}-u_n\ge3$, and
\[
 pQ\le(8p+8)K+4X+8p .
\]
Consequently $\sum(u_n^{-1/2}+(u_{n+1}-u_n-2)/u_n)\ge1/16$ over the charged
steps.
\end{theorem}

The integer inequality is checked in Lean and needs no minimality of $\tau$;
the $1/16$ corollary is an ordinary comparison over it.  The global energy
criterion --- that the parent hypotheses make the sequence eventually
Sylvester if and only if the record-energy series converges --- is not
supplied: the missing steps are odd-prime-power supply at every late index,
overlap selection across epochs, and the weighted comparison of the local
$1/16$ masses with a single series.

\begin{problem}[global epoch energy]\label{res:epochglobal}
Under the hypotheses of Problem~\ref{res:problem}, is the sequence eventually
Sylvester if and only if
\[
 \sum_{n\ \text{record}}\Bigl(\frac{\mathbf 1_{u_{n+1}-u_n\ge3}}{\sqrt{u_n}}
 +\frac{(u_{n+1}-u_n-2)_+}{u_n}\Bigr)<\infty\ ?
\]
\end{problem}

The inequality is the quantitative record-excess obligation that the first
batch left open: one protected prime power is reused across about $Q/8$
barriers, at the price of stopping at the scale $pQ/2$.  The square-root term
cannot be dropped, since a walk avoiding the multiples of a large prime with
jumps of $3$ only near those multiples has $\sum(u_{n+1}-u_n-2)/u_n=O(\log p/p)$
up to height $p^2/2$.

\begin{theorem}[saturated square transport]\label{res:saturated}
At every step, $u_{n+1}\mid h_ne_ne_{n+1}-v_n^2$.  In the normal form
$h_n=bc^2$, $a_n=bc\ell$, $v_n=bcs$ one has
$be_ne_{n+1}\equiv(bs)^2\pmod{u_{n+1}}$ with $b$, $s$ and $\ell$ units modulo
$u_{n+1}$.  Hence for every odd prime $p\mid u_{n+1}$ the Legendre symbols of
$e_ne_{n+1}$ and of $b$ agree, and a non-residue forces $h_n$ to be a
non-square, in particular $h_n\ge2$.
\end{theorem}

\begin{proof}
From $v_n-a_ne_n=(a_n-1)w_n$ and $h_ne_{n+1}=a_nv_n-(a_{n+1}-1)w_n$ one has
$a_ne_n\equiv v_n$ and $h_ne_{n+1}\equiv a_nv_n$ modulo $w_n$; multiplying
gives $w_n\mid v_n^2-h_ne_ne_{n+1}$ and $u_{n+1}\mid w_n$.  Substituting the
normal form and cancelling $bc^2$ as integers gives the refined congruence;
$v_{n+1}=b\ell s$ is coprime to $u_{n+1}$.
\end{proof}

All parts are checked in Lean.  The earlier transport of
Section~\ref{sec:open} worked modulo the part of $u_{n+1}$ coprime to
$G_{n+1}$; here the modulus is the whole numerator.  A pure-square payment
$h_n=c^2$ has $b=1$ and is invisible to the congruence: the family
$u=p^2k-1$, $a=pz$, $v=(a-1)u-1$ with $z\equiv1\ (k)$, $z\equiv0\ (k+1)$,
$p\nmid z$ has $e=e'=-1$ and $h=p^2$ with no defect.

\begin{lemma}[record-amplified error after cancellation]\label{res:amplified}
If $t$ is the first index after $s$ with $e_t<0$, then $u$ is nonincreasing on
$(s,t]$, so $u_t\le u_{s+1}$ and
$A_t\ge u_s/u_{s+1}=h_s/(1-e_s/u_s)$: a cancellation is visible at the next
negative error.
\end{lemma}

The visibility lemma is checked in Lean.  Pricing the deletion of a block
prime instead of protecting it is the mechanism that would reach every bound
$B$ where the parity barrier stops at $2$, but the global iteration ---
arbitrarily late CRT blocks, density-zero paid steps, and the transfer
$\limsup A_n<\infty$ if and only if eventual Sylvester recurrence --- is not
supplied.

\begin{problem}[global record-amplified criterion]\label{res:amplifiedglobal}
Assume the hypotheses of Problem~\ref{res:problem}.  Put $m_n=(-e_n)_+$,
$A_n=R_nm_n/u_n$ and $\delta_n=(a_n^2/a_{n+1}-1)_+$.  Is the sequence
eventually Sylvester if and only if $\limsup A_n<\infty$, if and only if
$\limsup R_n\delta_n<\infty$?  At the critical rate $\delta_n=O(1/n)$, with no
convergence assumption on $n\delta_n$, is it eventually Sylvester if and only
if $u_n=O(n)$, if and only if $(-e_n)_+=O(1)$?
\end{problem}

\begin{theorem}[two-modulus cut]\label{res:twomodulus}
Let $H\equiv5\pmod6$, and let $m\mid H$ and $\ell\mid H+2$ with $m,\ell>1$
dividing $v_{n+1}$.  Then no step has $u_n<H\le u_{n+1}$ with
$u_{n+1}-u_n\le4$.  Consequently, if late record-setting jumps are at most $4$
on a non-Sylvester orbit, then $\limsup\log G_n/\log u_n\ge4/3$ along records.
\end{theorem}

The cut is checked in Lean: the landings $H$ and $H+2$ die by primitivity, and
the landings $H+1$ and $H+3$ leave both endpoints even or both divisible by
$3$.  It is sharp, since $(u,v)=(33,5\cdot37\cdot41)$ with $a=231$ crosses
$H=35$ by a jump of $5$ with coprime endpoints.  The $4/3$ alternative charges
all but one of a block of pairwise coprime old moduli to $G_n$ by
Theorem~\ref{res:squarepayment}.

Three further ordinary results from the same batch are recorded in the working
note.  The tail gcd is sandwiched as $M_n\mid G_n\mid M_n^2$ against the LCM
overlap debt of Section~\ref{sec:lcmrecords}, so every lower bound on $G_n$
above is a lower bound $M_n\ge\sqrt{G_n}$.  No exact integer orbit agrees with a
rational multiple of a monic integer quadratic on more than a $3/4$ proportion
of any long interval, uniformly in the interval's position, by a Chebotarev
supply of prime obstructions.  The exceptional-discriminant negative-norm
calculation uses integer $k\ge1$; at $k=1/2$ both displayed norms are positive,
so that sign argument does not transfer unchanged to monic quadratics over
$\mathbb Q$.

No exact integer orbit agrees with a fixed rational cubic
$An(n+1)(n+2)+B$ on a density-one set of indices: each such profile has a
positive lower density of exceptions.  The reductions assume lower density
zero through gcd stabilisation, primitive shape, and the field-square
calculation; they do not yield a uniform density bound over all $A,B$.  The
two primitive words $2n(n+1)(n+2)\pm1$ are forbidden only at the stated
phases (plus starts $n\equiv0\pmod7$, minus starts $n\equiv1\pmod7$); disjoint
four-index blocks then give exceptional lower density at least $1/7$ on those
profiles.  The factor $1/(4\cdot7)=1/28$ is word-period accounting, not a uniform
constant, and an unrestricted ``no four consecutive agreements'' wording is
false: the plus profile admits a primitive four-step state at indices
$1,2,3,4$.  A middle zero at the forbidden phases forces $d_1^2\equiv2$ while
the coprime-update image is $\{0,2,5,6\}$.  The denominator-cleared identities
$(\alpha+2)Q(\alpha-1)-(\alpha-1)Q(\alpha)=3c$ and the resulting product identity
are checked in Lean over an arbitrary commutative ring; the divided form over
$\mathbb Q$ does not instantiate an irreducible number-field root.  Word
evaluation of the two mod-$7$ profiles is not orbit exclusion; exact three-step
transport of those words over $\mathbb F_7$ is empty.

The resulting cubic-rate statement is Theorem~\ref{res:cubicrate}
in the main text.

Every criterion in this section is a lower bound on a record, cancellation or
profile observable that the hypotheses of Problem~\ref{res:problem} leave free.
The missing implication is an upper bound on any one of them.  Eventual unit
record increments, under the stated prime-power supply, remain an exact
equivalence with the Sylvester endpoint.  Bounded $A_n$ and finite epoch energy
are proposed equivalences (Problems~\ref{res:amplifiedglobal}
and~\ref{res:epochglobal}).  The finite weighted record budget of
Section~\ref{sec:lcmrecords} remains the named equivalent open producer.  None
of these upper bounds has been derived from the hypotheses.



\section{Antecedents and statement conventions}\label{app:antecedents}

\noindent See~\cite[p.~64]{erdosgraham1980} and~\cite[p.~105]{erdos1988}.
Bloom's current catalogue record reproduces the displayed problem and labels
it open, while explicitly warning that the status is the website owner's
present assessment and may not reflect all relevant literature
\cite{erdosproblems}.  We therefore use the catalogue for numbering and
current reported status only; the original publications carry the
mathematical claims.  Below we index the same recurrence as
$a_{n+1}=a_n^{2}-a_n+1$, which is the shift the formal sources use; the two
forms are the same statement.  Write $\sylnext(a)=a^{2}-a+1$, the
\lword{Erdos243/ReciprocalTailRigidity.lean}{37}{sylvesterNext}{Sylvester
successor}.

\paragraph{Relation to prior work.}
Two published results reach the original problem under analytic hypotheses that
are not assumed below.  Their hypotheses are conditions on the sequence $(a_n)$
itself.  Duverney's is not comparable with the state-level hypotheses used
here; the Erd\H{o}s--Straus one, read through the dictionary of the next
paragraph, is implied by nonnegativity of $E$, the hypothesis of
Theorem~\ref{res:descent}, so that criterion already covers the descent proved
below (see~\cite[Cor.~4(1) and Prop.~1(1), pp.~14--15]{koizumi2025}).
Erd\H{o}s and Straus proved that if $\lim a_n/a_{n-1}^{2}=1$, the reciprocal
sum is rational, and $a_n$ does not satisfy the recurrence, then
\[
 \limsup_{n\to\infty}\ \frac{[a_1,\ldots,a_n]}{a_{n+1}}
 \left(\frac{a_{n+1}^{2}}{a_{n+2}}-1\right)>0,
\]
where $[a_1,\ldots,a_n]$ is the least common
multiple~\cite[Theorem~3, p.~132]{erdosstraus1964}.  This is the indexing in
the original theorem.  Erd\H{o}s's 1988 survey, followed by the current
catalogue summary, instead prints $a_n^2/a_{n+1}-1$ in the second factor while
retaining the same prefix-LCM quotient; that displayed summary is off by one
and is not followed here~\cite[p.~105]{erdos1988}\cite{erdosproblems}.
Koizumi records the convenient sufficient rate
$a_n^2/a_{n+1}=1+o(1/n)$ under which the Erd\H{o}s--Straus criterion settles
the problem~\cite[Remark~3, p.~16]{koizumi2025}.

Duverney proved a
conditional signed form of the problem itself: for signs
$\epsilon_n\in\{-1,1\}$, if
$\sum_{n\ge0}(a_{n+1}/a_n^{2}-1)$ converges, then the reciprocal sum with
numerators $\epsilon_n$ is rational if and only if
\[
 a_{n+1}=a_n^2-(\epsilon_{n+1}/\epsilon_n)a_n
             +\epsilon_{n+2}/\epsilon_{n+1}
\]
for all large $n$~\cite[Corollary~3.2, p.~287]{duverney2001}.  Absolute
convergence is a stronger sufficient specialisation, not the hypothesis
printed in that corollary.  The all-positive specialisation is the form
relevant here.

Badea's positive-term criterion is adjacent but different:
for positive rational terms $b_n/a_n$, eventual strict inequality
\[
 a_{n+1}>\frac{b_{n+1}}{b_n}a_n^2-\frac{b_{n+1}}{b_n}a_n+1
\]
forces irrationality, while rationality under the corresponding non-strict
inequality forces eventual equality~\cite[Theorem~A, p.~313; Cor.~2.2,
p.~316]{badea1993}.  These hypotheses are inequalities on successive denominators, not a bound on
the negative part of the cleared error.  Positive summands do not imply
nonnegative centred errors.  The comparison here concerns the quadratic-limit
regime, not recovery of either full published criterion.

Koizumi's work on doubly exponential sequences~\cite{koizumi2025} bears
directly on the state system used below, and supplies the coordinates in which
several of the statements of this note are most naturally read.  For a rational
$r=p/q$ with pseudo-greedy expansion $(a_n)$, remainders $(x_n)$ and gap
sequence $\varepsilon_n=x_n^{-1}+1-a_n$, his Lemma~4
~\cite[p.~11]{koizumi2025} produces integers $c_n>0$,
$d_n$ and $e_n$ with $x_n=c_n/d_n$ and $\varepsilon_n=e_n/c_n$, satisfying
\[
 a_n=\frac{d_n-e_n}{c_n}+1,\qquad
 c_{n+1}=c_n-e_n,\qquad
 d_{n+1}=a_nd_n ,
\]
and the proof of that lemma also gives $c_{n+1}=a_nc_n-d_n$.  Under the
dictionary $(C_n,D_n,E_n)=(c_n,d_n,e_n)$ these are exactly the objects of
Section~\ref{sec:state}: the first relation is the centring map $\ctr$ rewritten
as $e_n=d_n-(a_n-1)c_n$, the second is Proposition~\ref{res:update}, the third
is the denominator update, and the fourth is the tail update.  The lower-case
$e_n$ of Section~\ref{sec:state} is $-E_n$, so it is the negative of Koizumi's
$e_n$ at an index where the centred state is negative.

His Theorem~3
~\cite[pp.~12--13]{koizumi2025} shows
that his Conjecture~1, that for a positive rational $r$ whose gap sequence
satisfies $\varepsilon_n\to0$ one has $\varepsilon_n=0$ for all large $n$ ---
is equivalent to an affirmative answer to the question of Erd\H{o}s and Graham
recorded as his Question~1, which is Problem~\ref{res:problem} above.
Two of the implications proved below have prior art in these canonical
coordinates: his Lemma~2, that $\varepsilon_n=0$ forces $\varepsilon_{n+1}=0$,
is the absorption of Theorem~\ref{res:absorb}, and his Proposition~1(2), that
$\varepsilon_n\ge0$ for all large $n$ forces $\varepsilon_n=0$ for all large
$n$, is the descent of Theorem~\ref{res:descent}.  These appear in canonical
coordinates in~\cite[Lemma~2, p.~10; Prop.~1(2), p.~14]{koizumi2025}; no
question of priority or independence is
adjudicated here.

The growth hypothesis in Problem~\ref{res:problem} is calibrated by two facts
about Sylvester's sequence.  It satisfies $a_n\approx c_0^{2^{n}}$ with
$c_0=1.2640847\ldots$, and shifting it further produces sequences with
$a_n\approx C^{2^{n}}$ for arbitrarily large $C$ whose reciprocals still sum
to a rational number~\cite[arXiv v4, p.~2]{kovactao2024}.  The classical sufficient
condition for irrationality, $\lim_n a_n^{1/2^{n}}=\infty$, is therefore
sharp.  Kova\v{c} and Tao identify that condition as folklore
~\cite[arXiv v4, p.~2]{kovactao2024}; they attribute the sharpness observation to
Erd\H{o}s (1975).  A sequence with
$a_n/a_{n-1}^{2}\to1$ has $a_n^{1/2^{n}}$
convergent, so that criterion says nothing about the sequences of
Problem~\ref{res:problem}, and rationality is genuinely possible there.
Sylvester's
sequence is \texttt{A000058} in the OEIS and its shifts \texttt{A129871}.  For
the recent literature on irrationality of Ahmes series we refer to Kova\v{c}
and Tao~\cite{kovactao2024}, who resolve several problems of Erd\H{o}s and
Graham drawn from the same two sources cited above, and whose introduction
gives a sample of the intermediate work, including S\'andor (1984) and Badea
(1987).  They do not treat Problem~\#243; the rigidity conclusion asked for
there is not among their results.

The \emph{Formal Conjectures} collection contains a mathematically equivalent
unproved declaration, up to its zero-based indexing
~\cite{formalconjectures243}.  Its summand is $\Q$-valued, so Lean's
\texttt{Summable} hypothesis asserts the existence of a sum in $\Q$; the
finite indexing shift changes that sum only by a rational prefix.  The
declaration therefore does encode the rationality premise, but its proof is
\texttt{sorry}.  Its role here is statement-level prior art; it supplies no
proof authority, and the development below is independent of it.
The machine-checked results concern the state system: they exclude constant and
periodic negative magnitudes, the bounded-rise barrier of
Theorem~\ref{res:barrier}, and the two conditional endpoints
(Theorems~\ref{res:bounded} and~\ref{res:mass}), each stated with the exact
hypotheses that separate it from Problem~\#243.  No claim of priority is made
for anything below.

The hypothesis is asymptotic and the conclusion is exact, so the argument must
convert an analytic rate into an integer obstruction.  The conversion used here
is the classical one: clear denominators along the sequence, so that
rationality makes a tail integral, and then centre that integer at the value it
would take on Sylvester's sequence.  What remains is a single integer error
$E$, and the whole question is whether $E$ can avoid vanishing.


An author-posted preprint of Bado also studies bounded errors, LCM
compression and compensated mass~\cite{bado2026}.  Its Theorem~5.1 assumes
$|e_n|=O(1)$, rather than a bound only on the negative part.
This overlap calls for a dated comparison; it does not decide priority or
validate the preprint.


\section{Guide to the formal sources}\label{app:index}

Each linked phrase opens its Lean declaration at the pinned source
revision~\commitshort.  The state system, the exclusions, the barrier, and the
bounded-negative-part theorem are in \texttt{ReciprocalTailRigidity.lean}; the
finite-negative-mass theorem is in \texttt{SparseResetRecovery.lean}; and the
repair-entropy inequalities are in \texttt{RepairEntropy.lean}.  The residue
reduction of Appendix~\ref{app:residue} is in
\texttt{FiniteHorizonResidue.lean}.

The finite cores of Section~\ref{sec:records} are checked in the private
authoring root and are queued for the next pinned checkpoint, so they carry no
links at revision~\commitshort.
\begin{itemize}
\item \texttt{RecordIncrementBarrier.lean}: the per-step and telescoped square
  payment of Theorem~\ref{res:squarepayment}, the unit record-increment cut,
  and the composed consumer of Theorem~\ref{res:unitincrement}.
\item \texttt{SaturatedSquareTransport.lean}: the raw and normal-form square
  transport of Theorem~\ref{res:saturated}, the Legendre defect chain, and the
  record-amplified error lemma used in Lemma~\ref{res:amplified}.
\item \texttt{TwoModulusRecordCut.lean}: the cut of
  Theorem~\ref{res:twomodulus} and its first-crossing form.
\item \texttt{ProtectedEpochEnergy.lean}: the integer epoch inequality of
  Theorem~\ref{res:epoch}.
\item \texttt{CubicNeighbourIdentity.lean}: the denominator-cleared ring
  identities of Theorem~\ref{res:cubicrate}, the two primitive mod-$7$ words,
  and exact-transport impossibility of those words over $\mathbb F_7$.
\item \texttt{SignedDuverneyAlgebra.lean}: algebraic recurrences and cleared
  telescoping for the square-summable signed specialisation of Duverney.
\item \texttt{SparseResetRecovery.lean}: the scalar finite-mass lemma
  (Lemma~\ref{res:massscalar}) in addition to the vanishing form of
  Theorem~\ref{res:mass}.
\end{itemize}
The two-unit record-jump theorem of the first batch is in
\texttt{PrimitiveRecordBarrier.lean}.  Three distinctions are worth carrying into
the source.  The
identification of any of these modules with reciprocal tails is the exposition
of Section~\ref{sec:state} and is not a checked statement.  The periodic exclusion
assumes the regime $e_n<a_n$.  And the final Lean declaration lists strict
centring and normalised vanishing as hypotheses; the former follows eventually
from the latter with $K=1$, while the formal source derives neither from the
original analytic problem.  For the external bridge supplying them,
see Sections~\ref{sec:problem} and~\ref{sec:bounded}.

% The exhaustive source manifest is retained in the authored TeX for automated
% validation, and kept outside the printed note.
\iffalse

\begin{tabular}{@{}ll@{}}
informal statement & linked Lean source\\
Sylvester successor & \lrefx{Erdos243/ReciprocalTailRigidity.lean}{37}{sylvesterNext}\\
denominator update & \lrefx{Erdos243/ReciprocalTailRigidity.lean}{41}{nextDenState}\\
tail update & \lrefx{Erdos243/ReciprocalTailRigidity.lean}{45}{nextTailState}\\
centred state & \lrefx{Erdos243/ReciprocalTailRigidity.lean}{49}{centeredState}\\
Sylvester defect & \lrefx{Erdos243/ReciprocalTailRigidity.lean}{53}{sylvesterDefect}\\
update law & \lrefx{Erdos243/ReciprocalTailRigidity.lean}{57}{nextTailState_eq_sub_centered}\\
denominator equivariance & \lrefx{Erdos243/ReciprocalTailRigidity.lean}{65}{nextDenState_scale}\\
tail equivariance & \lrefx{Erdos243/ReciprocalTailRigidity.lean}{70}{nextTailState_scale}\\
centred equivariance & \lrefx{Erdos243/ReciprocalTailRigidity.lean}{75}{centeredState_scale}\\
normalised exclusion & \lrefx{Erdos243/ReciprocalTailRigidity.lean}{100}{no_normalizedConstantNegative_orbit}\\
scale-primitive exclusion & \lrefx{Erdos243/ReciprocalTailRigidity.lean}{147}{no_scalePrimitiveConstantNegative_orbit}\\
constant exclusion at every scale & \lrefx{Erdos243/ReciprocalTailRigidity.lean}{286}{no_constantNegative_orbit}\\
eventual constant exclusion & \lrefx{Erdos243/ReciprocalTailRigidity.lean}{350}{no_eventuallyConstantNegative_orbit}\\
persistence of a multiplier & \lrefx{Erdos243/ReciprocalTailRigidity.lean}{382}{base_dvd_denState_later}\\
one-step lock & \lrefx{Erdos243/ReciprocalTailRigidity.lean}{402}{divisorLocksTailSucc}\\
persistence of a lock & \lrefx{Erdos243/ReciprocalTailRigidity.lean}{429}{divisorLock_persists}\\
iterated period & \lrefx{Erdos243/ReciprocalTailRigidity.lean}{461}{periodic_add_mul}\\
iterated drift & \lrefx{Erdos243/ReciprocalTailRigidity.lean}{476}{phaseDrift_add_mul}\\
repeated-divisor transport & \lrefx{Erdos243/ReciprocalTailRigidity.lean}{494}{repeatedDivisor_forces_periodicCommonScale}\\
phase-primitive exclusion & \lrefx{Erdos243/ReciprocalTailRigidity.lean}{548}{no_phasePrimitivePeriodicNegative_orbit}\\
periodic exclusion at every scale & \lrefx{Erdos243/ReciprocalTailRigidity.lean}{690}{no_periodicNegative_orbit}\\
eventual periodic exclusion & \lrefx{Erdos243/ReciprocalTailRigidity.lean}{796}{no_eventuallyPeriodicNegative_orbit}\\
shifted consecutive multiples & \lrefx{Erdos243/ReciprocalTailRigidity.lean}{833}{exists_shifted_consecutiveMultiples}\\
bounded-rise barrier & \lrefx{Erdos243/ReciprocalTailRigidity.lean}{897}{no_boundedRise_of_tailAvoidance}\\
step coprimality & \lrefx{Erdos243/ReciprocalTailRigidity.lean}{983}{reducedStep_coprime_currentFactor}\\
pairwise coprimality & \lrefx{Erdos243/ReciprocalTailRigidity.lean}{1010}{reducedTail_pairwiseCoprime}\\
whole-modulus avoidance & \lrefx{Erdos243/ReciprocalTailRigidity.lean}{1024}{reducedTail_wholeModulusAvoidance}\\
reduced bounded-rise exclusion & \lrefx{Erdos243/ReciprocalTailRigidity.lean}{1035}{no_boundedRise_reducedTail}\\
eventual reduced exclusion & \lrefx{Erdos243/ReciprocalTailRigidity.lean}{1060}{no_eventuallyBoundedRise_reducedTail}\\
gcd monotonicity & \lrefx{Erdos243/ReciprocalTailRigidity.lean}{1424}{tailGcd_dvd_succ}\\
gcd at a negative index & \lrefx{Erdos243/ReciprocalTailRigidity.lean}{1581}{tailGcd_eq_gcd_negativeMagnitude}\\
chain stabilisation & \lrefx{Erdos243/ReciprocalTailRigidity.lean}{1599}{dvdChain_eventuallyConstant_of_cofinally_bounded}\\
gcd stabilisation & \lrefx{Erdos243/ReciprocalTailRigidity.lean}{1630}{tailGcd_eventuallyConstant_of_cofinally_boundedNegative}\\
tail-divergence form & \lrefx{Erdos243/ReciprocalTailRigidity.lean}{1655}{no_cofinallyBoundedNegative_of_tailTendsto}\\
divergence from normalised vanishing & \lrefx{Erdos243/ReciprocalTailRigidity.lean}{1733}{tailState_tendsto_atTop_of_nonzero_normalizedVanishes}\\
exclusion from normalised vanishing & \lrefx{Erdos243/ReciprocalTailRigidity.lean}{1748}{no_cofinallyBoundedNegative_of_normalizedVanishes}\\
defect identity & \lrefx{Erdos243/ReciprocalTailRigidity.lean}{1769}{sylvesterDefect_mul_nextTailState}\\
local rigidity step & \lrefx{Erdos243/ReciprocalTailRigidity.lean}{1781}{sylvesterNext_eq_of_centered_zero}\\
eventual Sylvester recurrence & \lrefx{Erdos243/ReciprocalTailRigidity.lean}{1799}{sylvesterNext_eventually_of_centered_zero}\\
eventual constancy & \lrefx{Erdos243/ReciprocalTailRigidity.lean}{1819}{antitone_nat_eventually_constant}\\
stabilisation of the nonnegative error & \lrefx{Erdos243/ReciprocalTailRigidity.lean}{1837}{centeredState_eventually_zero}\\
tail realisation & \lrefx{Erdos243/ReciprocalTailRigidity.lean}{1863}{natTail_eq_nextTailState}\\
denominator realisation & \lrefx{Erdos243/ReciprocalTailRigidity.lean}{1875}{natDen_eq_nextDenState}\\
signed update law & \lrefx{Erdos243/ReciprocalTailRigidity.lean}{1972}{natTail_eq_sub_centeredState}\\
absorption of a vanishing centred state & \lrefx{Erdos243/ReciprocalTailRigidity.lean}{2221}{centeredState_zero_absorbing}\\
contrapositive along a tail & \lrefx{Erdos243/ReciprocalTailRigidity.lean}{2248}{eventually_nonzero_of_zero_absorbing}\\
bounded-negative-part rigidity & \lrefx{Erdos243/ReciprocalTailRigidity.lean}{2265}{boundedNegativePart_eventually_zero}\\
eventual bounded-negative-part rigidity & \lrefx{Erdos243/ReciprocalTailRigidity.lean}{2352}{eventuallyBoundedNegativePart_eventually_zero}\\
repair-entropy conservation & \lrefx{Erdos243/RepairEntropy.lean}{48}{repairedFamily_recovery_energy_divisionFree}\\
independent-repair bound & \lrefx{Erdos243/RepairEntropy.lean}{78}{repaired_card_bound_of_independent}\\
fixed-length reset disappearance & \lrefx{Erdos243/RepairEntropy.lean}{107}{eventually_recoveryPayment_eq_one_of_fixedLength}\\
forced numerator & \lrefx{Erdos243/FiniteHorizonResidue.lean}{22}{forcedNumerator}\\
polynomial congruence & \lrefx{Erdos243/FiniteHorizonResidue.lean}{26}{forcedNumerator_modEq}\\
exact-division cancellation & \lrefx{Erdos243/FiniteHorizonResidue.lean}{41}{quotient_modEq_of_modEq_mul}\\
horizon modulus & \lrefx{Erdos243/FiniteHorizonResidue.lean}{53}{horizonModulus}\\
ascending-factorial form & \lrefx{Erdos243/FiniteHorizonResidue.lean}{59}{horizonModulus_eq_ascFactorial}\\
factorial value at the initial index & \lrefx{Erdos243/FiniteHorizonResidue.lean}{69}{horizonModulus_zero_eq_factorial}\\
survival predicate & \lrefx{Erdos243/FiniteHorizonResidue.lean}{75}{ForcedSurvives}\\
quotient transport & \lrefx{Erdos243/FiniteHorizonResidue.lean}{85}{forcedQuotient_modEq}\\
shrinking-modulus transport & \lrefx{Erdos243/FiniteHorizonResidue.lean}{97}{forcedSurvives_iff_of_modEq}\\
factorial residue reduction & \lrefx{Erdos243/FiniteHorizonResidue.lean}{134}{forcedSurvives_iff_of_modEq_factorial}\\
\end{tabular}
\fi

\section{A factorial residue reduction for forced orbits}\label{app:residue}

In the case $m=c=1$ of Section~\ref{sec:constant}, where the shape
equation~\eqref{eq:shape} reads $D_n+1=(a_n-1)(n+1)$, each multiplier is
determined by its predecessor; we call such an orbit \emph{forced}.  At index
$n$ the numerator of the next multiplier is
\[
 \fnum(n,a)=(n+1)a^{2}-(n+2)a+(n+3),
\]
the
\lword{Erdos243/FiniteHorizonResidue.lean}{22}{forcedNumerator}{forced
numerator}, and the divisor is $n+2$.  The orbit survives a step when that
division is exact, giving a survival predicate, the
\lword{Erdos243/FiniteHorizonResidue.lean}{75}{ForcedSurvives}{survival
predicate}.  Example~\ref{ex:shape} is the forced orbit from $a=3$ read this
way: $\fnum(0,3)=6$ is divisible by $2$ and gives $a_1=3$, while
$\fnum(1,3)=13$ is not divisible by $3$, so the orbit stops there.  Deciding
survival by iteration is expensive because the orbit
grows doubly exponentially, and it is unnecessary: survival over a finite
horizon depends on the initial value only through a factorial residue.

\begin{theorem}[factorial residue reduction]\label{res:residue}
For all $h$ and all integers $a\equiv b \pmod{(h+1)!}$, the orbit from $a$
survives $h$ forced updates exactly when the orbit from $b$ does.
\end{theorem}

\begin{proof}
Let $M(0,i)=1$ and $M(h+1,i)=(i+2)M(h,i+1)$, an ascending factorial with
$M(h,0)=(h+1)!$.  The numerator is a polynomial with integer coefficients, so
congruences transfer; reducing the modulus gives divisibility by $i+2$ for one
exactly when for the other, and cancelling that common factor from both values
and modulus leaves the inductive hypothesis at $M(h,i+1)$.
\end{proof}

\noindent Formalised as the
\lword{Erdos243/FiniteHorizonResidue.lean}{134}{forcedSurvives_iff_of_modEq_factorial}{factorial
residue reduction}, over the
\lword{Erdos243/FiniteHorizonResidue.lean}{97}{forcedSurvives_iff_of_modEq}{shrinking-modulus
transport}, the
\lword{Erdos243/FiniteHorizonResidue.lean}{26}{forcedNumerator_modEq}{polynomial
congruence}, and the
\lword{Erdos243/FiniteHorizonResidue.lean}{41}{quotient_modEq_of_modEq_mul}{exact-division
cancellation}; the modulus identifications are the
\lword{Erdos243/FiniteHorizonResidue.lean}{59}{horizonModulus_eq_ascFactorial}{ascending-factorial
form} and the
\lword{Erdos243/FiniteHorizonResidue.lean}{69}{horizonModulus_zero_eq_factorial}{factorial
value at the initial index}.

At $h=1$ the modulus is $2!=2$, and surviving one update means
$2\mid a^{2}-2a+3$, which holds exactly for odd $a$: for instance
$\fnum(0,3)=6$ but $\fnum(0,4)=11$.  So one step of survival is decided by the
parity of $a$ alone, which is Theorem~\ref{res:residue} at its smallest
nontrivial horizon.

The search this supported is superseded.  Running it over initial states below
$5000$ produced forced prefixes of length $17$ and no longer, which was
evidence and not a proof; Theorem~\ref{res:constant} now excludes the
constant-negative case outright, for every seed and at every scale.  The
reduction is retained because it is exact, and because the shrinking-modulus
technique transfers to any forced orbit whose step is a polynomial division.
